\section{Conclusions}
The result of this project is a reliable and easy-to-use application that automatically performs all the steps of the GRASP algorithm almost in real time. After reviewing the whole process from the initial ideas to the final testing, here are the main conclusions:
\smallskip\\\textbf{Main Objectives Achieved:} The main goal was to connect all the steps together in one automatic process, and this has been fully achieved. The application can now easily download data from AERONET and ACTRIS-EARLINET. Then, it runs the Matlab scripts to combine the data, executes the GRASP algorithm, and finally shows the results in a graph. All this automation saves time for the user and makes data analysis much faster.
\smallskip\\\textbf{Working on Different Systems and Handling Errors:} The decision to use C\# and the Avalonia UI package was very important. Thanks to this, the application works perfectly in the CALCULA system, which uses Linux. Also, the program is very strong at handling errors. Because different testing methods were used, such as unit testing and negative testing, the application can deal with problems, like missing data or changed file formats, without crashing. 
\smallskip\\\textbf{Better and Clearer Code:} At the beginning of the project, it was complicated to understand the old code because it was not organized and the file formats were outdated. However, this problem was solved by creating a more logical and structured program. By organizing the code following design patterns and good practices, it was also possible to create unit tests. This means that if someone else wants to work on this project in the future, it will be much easier to understand and improve.
\smallskip\\\textbf{Better User Experience:} The new application is designed specifically for the people who will use it. It has a clear interface that separates the main application settings from the specific project options, making everything simpler. Furthermore, there is an option to use a smaller interface designed for computer screens with a low resolution. This ensures that anyone can use the program, regardless of the computer they have.
\smallskip\\\textbf{Importance of Methodology and Requirements:} Finally, this project has shown how important it is to follow a clear methodology and have a well-defined list of requirements from the start. Because the initial requirements were not very clear, there were many unexpected problems and changes during the development. By eventually defining a strong plan and clear goals, it became much easier to solve problems, organize the code, and finish the application successfully. This proves that spending time on good planning and setting requirements is essential for any software project.
